\chapter{Simulation complexity versus decision complexity}
\label{ch:decision-complexity}

\section{Why simulator lower bounds are insufficient}
Showing that a circuit defeats one simulation strategy does not prove that its decision problem is outside $P$. A classical algorithm may compute the answer without reconstructing the quantum state. It is therefore insufficient to rule out only:

\begin{itemize}
\item full state-vector simulation,
\item one projector family,
\item bounded-width tensor-network contraction,
\item path sampling,
\item low-rank transfer matrices,
\item stabilizer decompositions,
\item any fixed catalogue of known simulators.
\end{itemize}

A proof of $P\neq BQP$ would need a lower bound against every deterministic polynomial-time classical algorithm for an explicit language.

\section{A representation-relative invariant}
For a declared classical representation class $\mathcal R$, define
\[
\operatorname{TC}^{\mathcal R}_\varepsilon(C)
=
\min\left\{
\operatorname{cost}(R):
R\in\mathcal R,
\ |R(C)-p_C|\le\varepsilon
\right\},
\]
where $p_C$ is the target acceptance probability. This permits a disciplined ladder:

\begin{enumerate}
\item exact polynomial-packet simulation for easy sectors;
\item lower bounds for local packet-query methods;
\item lower bounds for broader tensor, sampling, low-degree, and communication representations;
\item a representation-independent lower bound;
\item a separation of $P$ and $BQP$.
\end{enumerate}

Only the first step is completed here.

\section{Universal packetization and its consequence}
Consider the statement:

\begin{conjecture}[Universal polynomial packetization]
Every polynomial-time uniform quantum circuit family admits an efficiently discoverable exact or approximate packet representation of polynomial active size and polynomial transfer complexity, sufficient to compute its acceptance probability.
\end{conjecture}

If true, \cref{thm:poly-packet,thm:compressed-packet} would imply $BQP\subseteq P$, and hence $P=BQP$. The records do not establish this conjecture and explicitly identify global interference as an obstacle.

\section{The opposite route}
An explicit uniform family with a proved representation-independent classical lower bound and a polynomial quantum algorithm would establish $P\neq BQP$. The candidate succinct transfer-interference problem in the records illustrates the form of such a target, but the required lower bound is the unresolved theorem itself.

\begin{scopebox}
This monograph proves polynomial classical simulation for declared subclasses. It neither proves nor assumes a universal collapse of quantum computation.
\end{scopebox}

\section{Nonrelativizing issues}
Oracle separations show that query complexity can distinguish quantum and classical access. They also show why a black-box argument alone cannot settle ordinary $P$ versus $BQP$. A full separation would have to use explicit features such as uniform circuit generation, arithmetic structure, algebraic degree, communication constraints, or another nonrelativizing mechanism.
